\documentclass[10pt,a4paper]{amsart}
\usepackage[margin=19mm]{geometry}
\usepackage{amsmath,amssymb,mathtools}
\usepackage[scr=rsfs,cal=euler]{mathalfa}
\usepackage{xfrac}
\usepackage[unicode,hidelinks,bookmarks=false]{hyperref}
\newcommand{\ie}{i.e.\ }
\newcommand{\cf}{cf.\ }
\newcommand{\resp}{respectively\ }
\newcommand{\Subsection}{\subsection}
\providecommand{\nobreakdash}{\nobreak\hbox{-}}
\setlength{\emergencystretch}{2em}
\multlinegap0pt
\usepackage{bm}
\usepackage{bbm}
\usepackage{dsfont}
\usepackage{xspace}
\usepackage{cancel}
\usepackage{mathtools}
\usepackage{amsthm}
\makeatletter
\@ifundefined{theorem}{\newtheorem{theorem}{Theorem}}{}
\makeatother
\newcommand{\PP}{\mathbb{P}}
\title[Metropolis-Hastings Sampling of Phylogenetic Networks: Corre]{Metropolis-Hastings Sampling of Phylogenetic Networks: Correcting for Symmetries}
\author{Leo van Iersel, Remie Janssen, Mark Jones, Yukihiro Murakami, Christopher Reichling}
\date{Source: arXiv:2608.12430v1 (CC BY 4.0)}
\begin{document}
\maketitle

\noindent\emph{Source and licence.} Metropolis-Hastings Sampling of Phylogenetic Networks: Correcting for Symmetries. Version arXiv:2608.12430v1 (CC BY 4.0), distributed under the Creative Commons Attribution 4.0 International licence, as recorded in the arXiv licence metadata for this exact version and shown on the abstract page https://arxiv.org/abs/2608.12430v1. Full rights notice: licenses/arxiv-2608.12430-CC-BY-4.0.txt.\par\smallskip

\noindent\textit{Editorial scope.} Counted unit: the Theorem whose source label is \texttt{None} in the article (source0.tex lines 344--348, source label \texttt{None}), delivered together with the article's own complete original proof of it (source0.tex lines 349--383, one \texttt{proof} environment of 35 lines). No mathematics is written, repaired or re-derived: statement and proof are the article's own text line for line. Required context retained uncounted: none. Every definition, display and label of those lines is retained unchanged. Omitted from this package and not redistributed: 1--343, 384--1743. Nothing inside a retained excerpt is omitted or altered: every retained body is delivered exactly as the article's lines are written, so no omission receipt of any kind accompanies this package. Citations: the retained excerpts carry no citation command at all, so no bibliography entry of the article is transcribed and no reference is redistributed. Reference adaptation: none was needed and none was made; every reference inside the delivered unit resolves inside it, so no unresolved reference or citation is produced. Printed slips kept literal: no printed word, symbol, hypothesis or number of the counted statement or of its proof was altered. Layout and class: the article's own class is \texttt{article} with packages including fontenc, lmodern, inputenc, authblk, hyperref, natbib, doi, todonotes, amsmath, amsthm, amsfonts, mathtools, algorithm2e, cleveref; the wrapper is the lane's pinned \texttt{amsart} editorial wrapper at 10pt on A4, which restates the theorem-like environments the retained text needs and adds the article's own macro definitions that the retained text needs (\texttt{\textbackslash PP}), with extra wrapper packages (bm, bbm, dsfont, xspace, cancel, mathtools). The delivered numbering is the unit's own. Assets: no retained span contains a figure environment, a graphics-inclusion command, a PDF-page inclusion, a table or a TikZ picture; no image file is copied into this package, the package declares no asset dependency and no image file is redistributed with it. Licence: version arXiv:2608.12430v1 of this article is distributed under the Creative Commons Attribution 4.0 International licence, as recorded in the arXiv licence metadata for this exact version and shown on the abstract page https://arxiv.org/abs/2608.12430v1. A full rights notice is delivered at licenses/arxiv-2608.12430-CC-BY-4.0.txt. Characters: every retained excerpt is pure ASCII, so the delivered engine, XeLaTeX with its default Latin Modern text font, reports no missing character.
\begin{theorem}
    Let $\pi$ denote the stationary distribution of an irreducible aperiodic Markov chain $M=(X,T,P)$. Suppose we are given an equivalence relation $\sim$ partitioning the state space $X$ and the transitions $T$ of $M$ where the second equivalence respects the first, i.e. transitions can only be equivalent if they are between the same pair of aggregated states. Then the quotient Markov chain $M/\sim$ with transition probabilities
    \[P'_{\widetilde{m}} = \frac{1}{\pi(U)} \sum_{(m:u\to v)\in\widetilde{m}}\pi(u)P_m\]
    is irreducible and its unique stationary distribution $\widetilde\pi$ is compatible with $\pi$ in the sense that for every $V\in X/\sim$ we have $\widetilde\pi(V)=\pi(V)$.
\end{theorem}

\begin{proof}
    First, we check that the transition probabilities are actually transition probabilities, and the probabilities for all outgoing transitions sum to one for any state $U\in X/\sim$:
    \begin{align*}
    \sum_{\widetilde{m}: U\to \cdot} P'_{\widetilde{m}} 
        &= \sum_{\widetilde{m}: U\to \cdot} \frac{1}{\pi(U)} \sum_{(m:u\to v)\in\widetilde{m}}\pi(u) P_m \\
        &= \frac{1}{\pi(U)} \sum_{\widetilde{m}: U\to \cdot} \sum_{(m:u\to v)\in\widetilde{m}}\pi(u) P_m \\
        &= \frac{1}{\pi(U)} \sum_{u\in U} \sum_{m: u\to \cdot}\pi(u) P_m \\
        &= \frac{1}{\pi(U)} \sum_{u\in U} \pi(u) \sum_{m: u\to \cdot } P_m \\
        &= \frac{1}{\pi(U)} \sum_{u\in U} \pi(u) \\
        &= \frac{1}{\pi(U)} \pi(U)  \\
        &= 1
    \end{align*}

    Because the original Markov chain is aperiodic and irreducible, it is clear that the resulting Markov chain is aperiodic and irreducible as well. Hence, we only have to prove that the stationary distributions are compatible. For this, we simply calculate the stationary distribution for the quotient chain, which has transition probabilities $P'(V|U) := \sum_{\widetilde{m}: U \to V} P'_{\widetilde{m}}$ between two states $V,U\in X/\sim$.
    \begin{align*}
        \sum_{U\in X/\sim} \pi(U) P'(V|U) 
            &= \sum_{U\in X/\sim} \pi(U) \PP(V|U) \\
            &= \sum_{U\in X/\sim} \pi(U) \sum_{\widetilde{m}: U\to V} P'_{\widetilde{m}} \\
            &= \sum_{U\in X/\sim} \pi(U) \sum_{\widetilde{m}: U\to V} \frac{1}{\pi(U)} \sum_{(m:u\to v)\in\widetilde{m}}\pi(u) P_m \\
%            &= \sum_{U\in X/\sim} \pi(U)\frac{1}{\pi(U)} \sum_{\widetilde{m}: U\to V} \sum_{(m:u\to v)\in\widetilde{m}}\pi(u)p_m \\ 
            &= \sum_{U\in X/\sim} \sum_{\widetilde{m}: U\to V} \sum_{(m:u\to v)\in\widetilde{m}}\pi(u) P_m \\     
    \end{align*}
    Note that in the last summation, each transition to $V$ occurs exactly once, as both the states and the transitions are partitioned by $\sim$. Therefore, we can regroup the transitions by origin node and end node instead of equivalence class of transitions, we continue:
    \begin{align*}
        \sum_{U\in X/\sim} \pi(U) \PP(V|U) 
            % &= \sum_{U\in X/\sim} \sum_{u \in U} \sum_{v\in V} \sum_{m:u\to v}\pi(u)p_m \\
            % &= \sum_{v\in V} \sum_{U\in X/\sim} \sum_{u \in U} \sum_{m:u\to v}\pi(u)p_m \\
            &= \sum_{v\in V} \sum_{u \in X} \sum_{m:u\to v}\pi(u) P_m \\
            &= \sum_{v\in V} \sum_{u \in X} \pi(u) \sum_{m:u\to v} P_m \\
            &= \sum_{v\in V} \sum_{u \in X} \pi(u) P(v | u) \\
            &= \sum_{v\in V} \pi (v) \\
            &= \pi(V)
    \end{align*}

\end{proof}

\end{document}
